\scanpage{16}
\mspage{6}
% Archive foliation: 15, lower left.
On s'intéresse aux cas où les $f_{\mathcal H}^{*}$
ont des adjoints à g. $f_!^{\mathcal H}$ et à
droite $f_*^{\mathcal H}$, une condition sur $W$.

\uline{Remarque} Soit
\diagram{\begin{tikzcd}[row sep=small]
M\arrow[r,"\unsure{f|F}"]&M'\\
(W)&(W')
\end{tikzcd}}
tel que $f(W)\subset W'$. Alors
\[
\mathcal H_W(A)\longrightarrow\mathcal H_{W'}(A)
\]
\unsure{hom.}\inserted{bi}fonctoriel en $A$,
i.e. (\unsure{p.r. aux} $f_W^{*},f_{W'}^{*}$), ce qui
n'implique pas a priori qu'il commute
à iso près aux $f_!,f_*$.
Par ex\unsure{.} si $W,W'$ sont réduites aux iso, alors
cette commutation \deleted{\unsure{signifie}} est vérifiée
\inserted{pour $f_!$ (resp. $f_*$)}
ssi $f$ commute aux $\underset{\longrightarrow}{\lim}$
(resp. $\underset{\longleftarrow}{\lim}$).
Mais bien sûr
si $\mathcal H_W(A)\longrightarrow\mathcal H_{W'}(\unsure{A|A^{\circ}})$ est une équivalence
